\section{Basic notions}
\label{section-rings-basic}

\noindent
The following is a list of basic notions in commutative algebra. Some of these
notions are discussed in more detail in the text that follows and some are
defined in the list, but others are considered basic and will not be defined.
If you are not familiar with most of the italicized concepts, then we suggest
looking at an introductory text on algebra before continuing.

\begin{enumerate}
\item $R$ is a {\it ring},
\label{item-ring}
\item $x\in R$ is {\it nilpotent},
\label{item-ring-element-nilpotent}
\item $x\in R$ is a {\it zerodivisor},
\label{item-ring-element-zerodivisor}
\item $x\in R$ is a {\it unit},
\label{item-ring-element-unit}
\item $e \in R$ is an {\it idempotent},
\label{item-ring-element-idempotent}
\item an idempotent $e \in R$ is called {\it trivial} if $e = 1$ or $e = 0$,
\label{item-idempotent-trivial}
\item $\varphi : R_1 \to R_2$ is a {\it ring homomorphism},
\label{item-ring-homomorphism}
\item
\label{item-ring-homomorphism-finite-presentation}
$\varphi : R_1 \to R_2$ is {\it of finite presentation}, or
{\it $R_2$ is a finitely presented $R_1$-algebra},
see Definition \ref{definition-finite-type},
\item
\label{item-ring-homomorphism-finite-type}
$\varphi : R_1 \to R_2$ is {\it of finite type}, or
{\it $R_2$ is a finite type $R_1$-algebra},
see Definition \ref{definition-finite-type},
\item
\label{item-ring-homomorphism-finite}
$\varphi : R_1 \to R_2$ is {\it finite}, or
{\it $R_2$ is a finite $R_1$-algebra},
\item $R$ is a {\it (integral) domain},
\label{item-ring-domain}
\item $R$ is {\it reduced},
\label{item-ring-reduced}
\item $R$ is {\it Noetherian},
\label{item-ring-Noetherian}
\item $R$ is a {\it principal ideal domain} or a {\it PID},
\label{item-ring-PID}
\item $R$ is a {\it Euclidean domain},
\label{item-ring-Euclidean}
\item $R$ is a {\it unique factorization domain} or a {\it UFD},
\label{item-ring-UFD}
\item $R$ is a {\it discrete valuation ring} or a {\it dvr},
\label{item-ring-dvr}
\item $K$ is a {\it field},
\label{item-field}
\item $L/K$ is a {\it field extension},
\label{item-field-extension}
\item $L/K$ is an {\it algebraic field extension},
\label{item-field-extension-algebraic}
\item $\{t_i\}_{i\in I}$ is a {\it transcendence basis} for $L$ over $K$,
\label{item-transcendence-basis}
\item the {\it transcendence degree} $\text{trdeg}(L/K)$ of $L$ over $K$,
\label{item-transcendence-degree}
\item the field $k$ is {\it algebraically closed},
\label{item-algebraically-closed}
\item
\label{item-extend-into-algebraically-closed}
if $L/K$ is algebraic, and $\Omega/K$ an extension with $\Omega$
algebraically closed,
then there exists a ring map $L \to \Omega$ extending the map on $K$,
\item $I \subset R$ is an {\it ideal},
\label{item-ideal}
\item $I \subset R$ is {\it radical},
\label{item-ideal-radical}
\item if $I$ is an ideal then we have its {\it radical} $\sqrt{I}$,
\label{item-radical-ideal}
\item
\label{item-ideal-nilpotent}
$I \subset R$ is {\it nilpotent} means that $I^n = 0$ for
some $n \in \mathbf{N}$,
\item
\label{item-ideal-locally-nilpotent}
$I \subset R$ is {\it locally nilpotent} means that every
element of $I$ is nilpotent,
\item $\mathfrak p \subset R$ is a {\it prime ideal},
\label{item-prime-ideal}
\item
\label{item-prime-product-ideals}
if $\mathfrak p \subset R$ is prime and if $I, J \subset R$
are ideal, and if $IJ\subset \mathfrak p$, then
$I \subset \mathfrak p$ or $J \subset \mathfrak p$.
\item $\mathfrak m \subset R$ is a {\it maximal ideal},
\label{item-maximal-ideal}
\item any nonzero ring has a maximal ideal,
\label{item-exists-maximal-ideal}
\item
\label{item-jacobson-radical}
the {\it Jacobson radical} of $R$ is $\text{rad}(R) =
\bigcap_{\mathfrak m \subset R} \mathfrak m$ the intersection
of all the maximal ideals of $R$,
\item the ideal $(T)$ {\it generated} by a subset $T \subset R$,
\label{item-ideal-generated-by}
\item the {\it quotient ring} $R/I$,
\label{item-quotient-ring}
\item an ideal $I$ in the ring $R$ is prime if and only if $R/I$ is a domain,
\label{item-characterize-prime-ideal}
\item
\label{item-characterize-maximal-ideal}
an ideal $I$ in the ring $R$ is maximal if and only if the
ring $R/I$ is a field,
\item
\label{item-inverse-image-ideal}
if $\varphi : R_1 \to R_2$ is a ring homomorphism, and if
$I \subset R_2$ is an ideal, then $\varphi^{-1}(I)$ is an
ideal of $R_1$,
\item
\label{item-image-ideal}
if $\varphi : R_1 \to R_2$ is a ring homomorphism, and if
$I \subset R_1$ is an ideal, then $\varphi(I) \cdot R_2$ (sometimes
denoted $I \cdot R_2$, or $IR_2$) is the ideal of $R_2$ generated
by $\varphi(I)$,
\item
\label{item-inverse-image-prime}
if $\varphi : R_1 \to R_2$ is a ring homomorphism, and if
$\mathfrak p \subset R_2$ is a prime ideal, then
$\varphi^{-1}(\mathfrak p)$ is a prime ideal of $R_1$,
\item $M$ is an {\it $R$-module},
\label{item-module}
\item
\label{item-annihilator}
for $m \in M$ the {\it annihilator}
$I = \{f \in R \mid fm = 0\}$ of $m$ in $R$,
\item $N \subset M$ is an {\it $R$-submodule},
\label{item-submodule}
\item $M$ is a {\it Noetherian $R$-module},
\label{item-Noetherian-module}
\item $M$ is a {\it finite $R$-module},
\label{item-finite-module}
\item $M$ is a {\it finitely generated $R$-module},
\label{item-finitely-generated-module}
\item $M$ is a {\it finitely presented $R$-module},
\label{item-finitely-presented-module}
\item $M$ is a {\it free $R$-module},
\label{item-free-module}
\item
\label{item-extension-free}
if $0 \to K \to L \to M \to 0$ is a short exact sequence
of $R$-modules and $K$, $M$ are free, then $L$ is free,
\item if $N \subset M \subset L$ are $R$-modules, then $L/M = (L/N)/(M/N)$,
\label{item-isomorphism-theorem}
\item $S$ is a {\it multiplicative subset of $R$},
\label{item-multiplicative-subset}
\item the {\it localization} $R \to S^{-1}R$ of $R$,
\label{item-localization-ring}
\item
\label{item-localization-zero}
if $R$ is a ring and $S$ is a multiplicative subset
of $R$ then $S^{-1}R$ is the zero ring if and only if $S$ contains $0$,
\item
\label{item-localize-nonzerodivisors}
if $R$ is a ring and if the multiplicative subset $S$
consists completely of nonzerodivisors, then $R \to S^{-1}R$
is injective,
\item if $\varphi : R_1 \to R_2$ is a ring homomorphism, and
$S$ is a multiplicative subset of $R_1$, then $\varphi(S)$ is
a multiplicative subset of $R_2$,
\item
\label{item-products-multiplicative-subsets}
if $S$, $S'$ are multiplicative subsets of $R$,
and if $SS'$ denotes the set of products $SS' =
\{r \in R \mid \exists s\in S, \exists s' \in S', r = ss'\}$
then $SS'$ is a multiplicative subset of $R$,
\item
\label{item-localization-localization}
if $S$, $S'$ are multiplicative subsets of $R$,
and if $\overline{S}$ denotes the image of $S$ in $(S')^{-1}R$,
then $(SS')^{-1}R = \overline{S}^{-1}((S')^{-1}R)$,
\item the {\it localization} $S^{-1}M$ of the $R$-module $M$,
\label{item-localization-module}
\item
\label{item-localization-exact}
the functor $M \mapsto S^{-1}M$ preserves injective maps,
surjective maps, and exactness,
\item
\label{item-localization-localization-module}
if $S$, $S'$ are multiplicative subsets of $R$,
and if $M$ is an $R$-module, then
$(SS')^{-1}M = S^{-1}((S')^{-1}M)$,
\item
\label{item-localize-ideal}
if $R$ is a ring, $I$ an ideal of $R$, and $S$ a multiplicative
subset of $R$, then $S^{-1}I$ is an ideal of $S^{-1}R$, and we have
$S^{-1}R/S^{-1}I = \overline{S}^{-1}(R/I)$, where $\overline{S}$
is the image of $S$ in $R/I$,
\item
\label{item-ideal-in-localization}
if $R$ is a ring, and $S$ a multiplicative
subset of $R$, then any ideal $I'$ of $S^{-1}R$ is
of the form $S^{-1}I$, where one can take $I$ to be
the inverse image of $I'$ in $R$,
\item
\label{item-submodule-in-localization}
if $R$ is a ring, $M$ an $R$-module, and $S$ a multiplicative
subset of $R$, then any submodule $N'$ of $S^{-1}M$ is of the form
$S^{-1}N$ for some submodule $N \subset M$, where
one can take $N$ to be the inverse image of $N'$ in $M$,
\item if $S = \{1, f, f^2, \ldots\}$ then $R_f = S^{-1}R$ and $M_f = S^{-1}M$,
\label{item-localize-f}
\item
\label{item-localize-p}
if $S = R \setminus \mathfrak p = \{x\in R \mid x\not\in \mathfrak p\}$
for some prime ideal $\mathfrak p$,
then it is customary to denote $R_{\mathfrak p} = S^{-1}R$
and $M_{\mathfrak p} = S^{-1}M$,
\item a {\it local ring} is a ring with exactly one maximal ideal,
\label{item-local-ring}
\item a {\it semi-local ring} is a ring with finitely many maximal ideals,
\label{item-semi-local-ring}
\item
\label{item-localize-p-local-ring}
if $\mathfrak p$ is a prime in $R$, then $R_{\mathfrak p}$ is
a local ring with maximal ideal $\mathfrak p R_{\mathfrak p}$,
\item
\label{item-residue-field}
the {\it residue field}, denoted $\kappa(\mathfrak p)$,
of the prime $\mathfrak p$ in the ring $R$ is the
field of fractions of the domain $R/\mathfrak p$;
it is equal to $R_\mathfrak p/\mathfrak pR_\mathfrak p
= (R \setminus \mathfrak p)^{-1}R/\mathfrak p$,
\item given $R$ and $M_1$, $M_2$ the {\it tensor product} $M_1 \otimes_R M_2$,
\label{item-tensor-product}
\item
\label{item-cauchy-binet}
given matrices $A$ and $B$ in a ring $R$ of sizes $m \times n$ and
$n \times m$ we have $\det(AB) = \sum \det(A_S)\det({}_SB)$ in $R$ where
the sum is over subsets $S \subset \{1, \ldots, n\}$ of size $m$
and $A_S$ is the $m \times m$ submatrix of $A$ with columns
corresponding to $S$ and ${}_SB$ is the $m \times m$ submatrix of $B$
with rows corresponding to $S$,
\item etc.
\end{enumerate}




